\documentclass[master,paper]{pmk-graduate}
% Вид документа
% bachelor - Бакалавриат ВКР
% master - Магистратура ВКР
% term - Курсовая работа
% essay - Реферат
%
% assignment - Задание на работу
% paper - Текст работы
%
% thesis - Текст автореферата

\usepackage{amsmath}
\usepackage{multirow}
\usepackage{diagbox}
\usepackage{tabularx}
\usepackage{booktabs}
\usepackage{listings}
\usepackage{xcolor} % Для подсветки синтаксиса

\addbibresource{bondal.bib}

\begin{document}
% Учебный год
\year{2025-2026}

% Приказ: дата, номер
%\order{02.12.2025}{222-С}

% Дисциплина для курсовой работы и реферата
% \subject{Дискретная математика}


% Студент
\student{Бондал Андрей Алексеевич
\brief А.\,А.\,Бондал
\group М-22
\english Bondal Andrey Alekseevich
}

% Научный руководитель
\superviser{Дудаков Сергей Михайлович
\brief С.\,М.\,Дудаков
\position декан факультета ПМиК
\degree д.\,ф.-м.\,н. % Учёная степень. При отсутствии убрать строку
\rank профессор % Учёное звание. При отсутствии убрать строку
}

% Образовательная программа. См. перечень ниже
\program{mista}
% Бакалавриат:
% ics - Информатика и компьютерные науки
% swi - Инженерия программного обеспечения
% swiai - Программная инженерия в искусственном интеллекте
% aida - Искусственный интеллект и анализ данных
% mm - Математическое моделирование
% sa - Системный анализ
% icpim - Инженерное компьютерное проектирование и индустриальная математика
% acse - Прикладная информатика в экономике
% acsm - Прикладная информатика в мехатронике
% icmrs - Интеллектуальное управление в мехатронных и робототехнических системах
% Магистратура:
% macsae - Прикладная информатика в аналитической экономике
% mista - Интеллектуальные системы. Теория и приложения
% msp - Системное программирование
% msa - Системный анализ
% mitcd - Информационные технологии в управлении и принятии решений


% Заголовок ВКР
\title{Интеллектуальные методы коррекции базы данных
    \english Intelligent database correction methods}

% Цель ВКР
\goal{Разработка системы для коррекции ошибок в базах данных, вызванных нарушением ограничений целостности.}

% Задачи для разработки в ВКР
\begin{tasks}
\task{Описать выбранный метод, его теоретические основы и возможности использования.
\result %Указывается только для самой работы. В задании можно не писать
В рамках данного исследования был описан и разработан метод, сочетающий в себе принципы генетического алгоритма и глубокое понимание структуры и особенностей \texttt{SQL}-запросов. Ключевым аспектом является адаптация генетических операторов (селекция, кроссовер, мутация) для работы с \texttt{SQL}-запросами, а также разработка специализированной фитнес-функции, учитывающей не только синтаксическую корректность запросов, но и соответствие схеме БД, соблюдение ограничений целостности, а также минимизацию изменений данных. Метод включает в себя следующие этапы:
\begin{enumerate}
	\item анализ \texttt{SQL}-запроса и выявление нарушений (ограничений целостности);
	\item генерация начальной популяции \texttt{SQL}-запросов, представляющих собой потенциальные исправления ошибки;
	\item оценка каждого запроса в популяции с использованием фитнес-функции, учитывающей несколько критериев;
	\item применение генетических операторов (селекция, кроссовер, мутация) для создания нового поколения запросов;
	\item повторение шагов <<в>> и <<г>> до достижения критерия остановки;
	\item выбор наилучшего \texttt{SQL}-запроса из популяции, который представляет собой оптимальное исправление ошибки.
\end{enumerate}
В отличие от существующих подходов, предложенный метод позволяет учитывать различные типы ограничений и находить оптимальные решения даже в сложных сценариях.
}    
\task{Реализовать предложенный метод в контексте коррекции базы данных.
\result %Указывается только для самой работы. В задании можно не писать
Для практической проверки эффективности предложенного метода была разработана система на языке \texttt{Java}, которая включает в себя следую\-щие ключевые компоненты:
\begin{enumerate}
	\item парсер \texttt{SQL}, который обеспечивает анализ входных \texttt{SQL}-запросов и выявление ошибок;
	\item генератор \texttt{SQL}, который генерирует новые \texttt{SQL}-запросы на основе генетических операторов и обеспечивает синтаксическую корректность запросов, а так же соблюдение ограничений схемы базы данных;
	\item генетические операторы, а именно операторы турнирной селекции, одноклеточного кроссовера и мутации, адаптированные для работы с \texttt{SQL}-запросами;
	\item фитнес-функция, которая позволяет оценивать качество \texttt{SQL}-запро\-сов на основе нескольких критериев (соблюдение ограничений целостности, соответствие схеме базы данных, минимизация изменений данных и т.д.).
\end{enumerate}
Разработанная система представляет собой гибкий и расширяемый инструмент для коррекции ошибок в базах данных, а архитектура позволяет легко добавлять новые генетические операторы, критерии оценки и СУБД.
}    
\task{Оценить эффективность метода на практике и проанализировать полученные результаты.
\result %Указывается только для самой работы. В задании можно не писать
Для оценки эффективности разработанного метода и определения оптимальных значений параметров генетического алгоритма была проведена серия экспериментов на тестовых данных. Тестовые данные включали в себя набор \texttt{SQL}-запросов с различными типами ошибок и тестовую базу данных с заранее определенной структурой и ограничениями. В ходе экспериментов варьировались параметры генетического алгоритма (размер популяции, вероятность кроссовера, вероятность мутации и т.д.) и оценивалось качество полученных решений (количество успешно исправленных ошибок, время работы алгоритма и т.д.). Результаты экспериментов показали, что предложенный метод позволяет эффективно корректировать ошибки в базах данных. Были определены следующие оптимальные значения параметров генетического алгоритма:
\begin{itemize}
	\item размер популяции: 50;
	\item вероятность кроссовера: 0.8;
	\item вероятность мутации: 0.1;
	\item количество поколений: 100.
\end{itemize}
Полученные результаты подтверждают работоспособность и эффективность предложенного подхода. Оптимальные значения параметров генетического алгоритма позволяют достичь наилучшего баланса между качеством решений и временем работы алгоритма.}    
\end{tasks}

% Перечень иллюстративного материала
%\begin{images}
%\image{1}    
%\image{2}    
%\end{images}

% Консультанты
%\begin{advisers}
%\adviser{Сидоров Сидор Сидорович, проектирование схемы БД}    
%\end{advisers}

% Команда для генерации задания на ВКР
%\makeassignment
%\end{document}

% Далее идёт текст самой работы

% Английская аннотация
\begin{englishabstract}
This research introduces a novel approach to address the persistent challenge of data integrity violations in modern databases by leveraging a genetic algorithm-based error correction system. The methodology commences with a comprehensive analysis of database schemas and pre-defined integrity constraints. Subsequently, diagnostic error messages generated by the database system are intelligently utilized to construct a pool of potential \texttt{SQL} queries designed to rectify the identified inconsistencies and restore data integrity.

The core of the system lies in an evolutionary process, driven by a carefully crafted fitness function. This function simultaneously optimizes several key criteria: adherence to the database schema, satisfaction of specified integrity constraints, preservation of existing data to minimize unintended side effects, and minimization of placeholder values introduced during the repair process, ensuring data quality and accuracy.

While acknowledging the inherent computational intensity required to locate globally optimal repairs within complex database systems, the proposed algorithm demonstrates its capacity to identify high-quality, data-consistent solutions through iterative refinement and a strategic exploration of the vast search space. This targeted exploration allows for efficient convergence towards viable repair strategies, balancing computational cost with solution effectiveness. Preliminary experimental results showcase the algorithm’s robustness and effectiveness in resolving diverse types of constraint violation scenarios, highlighting its potential for practical application in real-world database management systems. Further research will focus on optimizing the algorithm’s performance and expanding its applicability to a wider range of data integrity challenges.
\end{englishabstract}

\maketitlepage
\introduction
\actuality
Современные информационные системы не могут обойтись без баз данных. Они играют важную роль в хранении и обработке критически важной информации для различных отраслей -- от бизнеса до науки.

Задача поддержания целостности данных -- это одно из ключевых направлений работы с базами данных, ведь любые ошибки в них могут привести к сбоям в работе системы и неверным аналитическим выводам. Корень проблемы может возникать по разным причинам: человеческий фактор, недоработки программного обеспечения или даже злонамеренные действия.

Традиционные способы гарантии согласованности данных, такие как использование ограничений целостности, помогают выявлять и предотвращать появление ошибок. Однако часто они недостаточны для их устранения. Ручная корректировка исправления ошибок является крайне трудоемким и дорогостоящим процессом. Требуется наличие высококвалифицированных специалистов. Именно поэтому разработка интеллектуальных методов исправления ошибок в базах данных приобретает большое значение.

\review
Для решения проблемы исправления ошибок в базах данных в этом иссле\-довании предлагается использовать генетический алгоритм.

Основные принципы и теоретические аспекты работы генетических алгоритмов подробно описаны в учебнике \cite{bib1}, который был основой для изучения и применения этого метода. Математический аппарат генетических алгоритмов, включая способы представления генотипа и функции пригодности, был изучен с помощью статьи на сайте Loginom \cite{bib2}, что позволило углубить понимание теоретических основ. Статья на сайте Интуит \cite{bib3}, посвященная генетическому программированию, позволила ознакомиться с альтернативными подходами к решению задач автоматизации и оптимизации, хотя и не являет\-ся основным методом в данном исследовании.

В то время как для понимания теоретических аспектов целостности данных, была изучена статья \cite{bib4}, которая рассматривает сложность обновлений дедуктивных баз данных с фиксированными ограничениями целостности, но задача исправления ошибок, требующая создания и проверки гипотез о возможных коррекциях остается малоизученной.

Для разработки программной реализации предложенного алгоритма использовалось руководство по языку программирования \texttt{Java} \cite{bib5}, обеспечи\-вающее необходимую информацию для эффективного использования возможностей языка и стандартных библиотек, особенно в части работы с базами данных (\texttt{JDBC}). Практические шаги реализации генетического алгоритма были изучены с помощью статьи \cite{bib6}, что позволило ускорить процесс разработки. Для создания удобного интерфейса взаимодейст\-вия с генетическим алгоритмом и базой данных, была изучена статья \cite{bib7}, посвященная \texttt{Java Swing}.

Официальная документация \texttt{PostgreSQL} \cite{bib8} предоставила основной источник информации о работе с данной СУБД, необходимой для хранения данных, выполнения запросов и управления транзакциями.

Таким образом, данная работа нацелена на заполнение пробела в научных исследованиях и разработку практичного подхода к устранению ошибок, возникающих из-за нарушения целостности данных, путем применения генетического алгоритма, что предполагает использование знаний, полученных из изученных источников, для решения поставленной задачи.

\goalandtasks

\tools
В работе для исправления ошибок в базах данных использовался генетический алгоритм. Он генерирует различные варианты исправления \texttt{SQL}-запросов и оценивает, насколько хорошо они соответствуют требованиям к базе данных. Чтобы сгенерированные запросы работали правильно, применялся анализ структуры базы данных. Всё это было реализовано на языке \texttt{Java} с использованием библиотеки \texttt{JDBC} для взаимодействия с базой данных и \texttt{Swing} для создания графического интерфейса.

\results

\originality
Новизна работы заключается в разработке и применении специализированной фитнес-функции, учитывающей не только соответствие \texttt{SQL}-запросов ограничениям целостности и схеме базы данных, но и эффективность исправления ошибок и сохранение данных, а также в адаптации генетических операторов для работы с \texttt{SQL}-запросами. Спецификой разработанной прог\-раммы является возможность динамической настройки параметров генетического алгоритма и учет различных типов ошибок при оценке решений.

\approbation
Основные итоги были представлены на студенческой конференции, проводимой факультетом прикладной математики и кибернетики Тверского государственного университета в 2025 г. Результаты работы были обсуждены и одобрены на кафедре информатики.

\structure
Работа будет организована следующим образом:

Введение содержит обоснование актуальности темы исследования, обзор литературы, цели и задачи работы, описание используемых методов и средств, перечень основных результатов, представление научной новизны, сведения об апробации, а также описание структуры работы.

В главе \ref{chap1} излагаются теоретические основы генетических алгоритмов (ГА), необходимые для понимания и реализации разработанного метода коррекции ошибок.

В параграфе \ref{sec11} даются определения ключевых терминов и понятий ГА, используемых в работе, тем самым позволяя избежать неоднозначности и обеспечить единое понимание терминологии.

Параграф \ref{sec12} содержит подробное описание этапов работы генетического алгоритма, от инициализации популяции и до критериев остановки.

В параграфе \ref{sec13} описывается способ представления решений задачи коррекции БД в виде хромосом, пригодных для обработки генетическим алгоритмом, и другие виды кодирования.

Глава \ref{chap2} посвящена детальному описанию разработанного ГА для решения задачи коррекции базы данных, включая особенности его реализации и применяемые подходы.

В параграфе \ref{sec21} дается общее описание разработанной системы, реализующей ГА коррекции БД, описывается архитектура, компоненты системы и их взаимодействие.

Параграф \ref{sec22} содержит описание фитнес-функции, которая используется для оценки качества решений (хромосом).

Выбранные и реализованные генетические операторы и параметры ГА описываются в параграфе \ref{sec23}

Параграф \ref{sec24} демонстрирует пользовательский интерфейс разработанной программы с особенностями его использования.

В главе \ref{chap3} представляются итоги экспериментальных исследований разработанного алгоритма.

Так в параграфе \ref{sec31} описываются данные, которые используются для проведения экспериментальных исследований.

Параграф \ref{sec32} описывает процесс проведения тестов, примеров и анализа полученных результатов, так же в нём предоставляются выводы об эффективности разработанного алгоритма.

В заключении представлены ключевые результаты исследования, а также обозначены перспективные направления для дальнейшей работы.

Раздел «Список литературы» включает все использованные в работе источники.

В приложении \ref{chap00} находится код разработанной системы.

\chapter{Методология генетических алгоритмов}
\label{chap1} 
\section{Основные понятия}
\label{sec11} 

Генетический алгоритм (ГА) – это эвристический метод поиска и оптимизации, основанный на принципах естественного отбора и генетики. Как указано в пособии \cite{bib1}, ГА моделирует процесс эволюции биологических видов, где наиболее приспособленные особи имеют больше шансов на выживание и продолжение рода, передавая свои генетические признаки потомству. Этот процесс повторяется из поколения в поколение, приводя к улучшению характеристик популяции.

В основе ГА лежит идея, что решение сложной задачи можно найти, <<выращивая>> его из набора случайных решений (популяции) с помощью итеративного процесса, имитирующего естественный отбор. Лучшие решения (более приспособленные особи) имеют больше шансов на <<размножение>> (генерацию новых решений на их основе), а худшие решения отбраковываются. Операторы кроссовера и мутации имитируют генетические процессы, обеспечивая разнообразие популяции и возможность появления новых, более эффективных решений.

Ключевые термины:
\begin{enumerate}
	\item \textit{популяция} -- набор решений (хромосом), представляющих собой возможные варианты решения задачи. Размер популяции – важный параметр ГА, влияющий на его эффективность;
	\item \textit{хромосома} -- представление решения задачи в виде структуры данных (обычно бинарной строки, но могут быть и другие варианты). В рамках коррекции баз данных, хромосома может представлять собой \texttt{SQL}-запрос;
	\item \textit{ген} -- компонент хромосомы, представляющий собой отдельную характеристику решения. Например, в \texttt{SQL}-запросе ген может представлять собой значение атрибута, оператор сравнения или условие;
	\item \textit{генотип} -- совокупность генов хромосомы, т.е. полное описание решения;
	\item \textit{фенотип} -- фактическое решение, которое получается из генотипа. В нашем случае, это \texttt{SQL}-запрос, который можно выполнить в базе данных;
	\item \textit{фитнес-функция} -- функция, оценивающая качество решения (хромосомы) и возвращающая значение, отражающее его приспособленность. Чем оно выше, тем лучше решение. В задаче коррекции баз данных (БД), фитнес-функция должна учитывать такие факторы, как соблюдение ограничений целостности, соответствие схеме БД, эффективность исправления и сохранение данных;
	\item \textit{поколение} -- одна итерация работы генетического алгоритма, включаю\-щая оценку приспособленности, селекцию, кроссовер и мутацию.
\end{enumerate}

Понимание этих понятий необходимо для успешного применения ГА к задаче коррекции баз данных.

\section{Этапы работы генетического алгоритма}
\label{sec12} 

Генетический алгоритм представляет собой итеративный процесс, состоящий из нескольких ключевых этапов, которые последовательно применяются к популяции решений, а именно:

\begin{enumerate}
	\item \textit{инициализация популяции};
	\item \textit{оценка приспособленности};
	\item \textit{селекция (отбор)};
	\item \textit{кроссовер (скрещивание)};
	\item \textit{мутация};
	\item \textit{замена популяции};
	\item \textit{критерии остановки}.
\end{enumerate}

В начале работы генетического алгоритма формируется популяция случайных решений (хромосом). Как упоминалось ранее, размер популяции играет важную роль в эффективности работы алгоритма. Начальная популяция должна быть достаточно разнообразной для обеспечения широкого охвата пространства поиска решений. В контексте коррекции баз данных на этапе инициализации генерируются случайные \texttt{SQL}-запросы - они могут быть как синтаксически верными, так и содержать ошибки. Важно учитывать схему БД и ограничения целостности при создании запросов для избежания генерации нерабочих решений.

На этапе \textit{оценки приспособленности} для каждой хромосомы в популяции вычисляется значение фитнеса, которое отражает её пригодность к решению поставленных задач. Фитнес играет важную роль в ГА, поскольку опреде\-ляет те варианты решений, которые считаются <<наилучшими>>, и их отбор для пос\-ледующего размножения природным путем. В задаче коррекции баз данных фитнес должен учитывать много факторов - такие как соблюдение ограничений целостности данных, соответствие структуре БД, эффективность исправ\-лений и сохранение информации.

На следующем этапе отбираются наиболее приспособленные хромосомы из популяции для участия в процессе размножения. Существует множество методов \textit{селекции}, каждый из которых имеет свои преимущества и недостатки. В пособии рассматривается два вида: \textit{турнирная} и \textit{рулеточная селекция}. В первом случае происходит случайный выбор хромосом из популяции и отбор лучшей из них. А во втором -- пропорциональный отбор хромосом, где вероятность выбора хромосомы пропорциональна ее значению фитнес-функции.

\textit{Кроссовер} -- этап, на котором создаются новые хромосомы путем комбинирования генетического материала двух родительских хромосом. Он имити\-рует процесс полового размножения в природе. Существует множество методов \textit{кроссовера}, таких как: \textit{одноточечный}, \textit{многоточечный} и \textit{равномерный кроссовер}. Рассмотрим каждый из них:
\begin{itemize}
\item при \textit{одноклеточном кроссовере} выбирается случайная точка разрыва в обеих родительских хромосомах. Новые хромосомы создаются путем обмена частями хромосом до и после точки разрыва;
\item в \textit{многоклеточном} методе выбирается несколько точек разделения в родительских хромосомах. Новые же хромосомы создаются как при одноклеточном способе;
\item при \textit{равномерном кроссовере} каждый ген в новой хромосоме выбирает\-ся случайным образом у одного из родителей.
\end{itemize}

\textit{Мутация} -- этап, на котором в хромосомы вносятся случайные изменения. Он имитирует процесс мутаций в природе и обеспечивает генетическое разнообразие популяции. Вероятность мутации обычно невелика, чтобы избежать разрушения хороших решений. Существуют различные типы мутаций, такие как: \textit{битовая мутация}, \textit{инверсия} и \textit{перестановка}. При \textit{битовой мутации} каждый бит в бинарной строке инвертируется (допустимо при бинарном представлении данных), в \textit{инверсии} -- выбирается участок хромосомы и инвертируется порядок генов на участке, а при \textit{перестановке} случайно выбираются два гена в хромосоме и меняются местами.

На этапе \textit{замены популяции} старая популяция заменяется новой, состоящей из отобранных родительских хромосом и новых, созданных в результате кроссовера и мутации. Существуют различные стратегии замены популяции, такие как полная замена (все хромосомы заменяются новыми) и частичная (только часть хромосом заменяется новыми).

Так как генетический алгоритм является итеративным процессом, то он повторяется до тех пор, пока не будет достигнут один из \textit{критериев остановки}. Наиболее распространенные критерии:
\begin{itemize}
\item достижение максимального количества поколений;
\item нахождение решения, удовлетворяющего заданному критерию качества;
\item отсутствие улучшений в течение заданного количества поколений.
\end{itemize}

Каждый этап играет важную роль в работе алгоритма и влияет на его эффективность. Правильный выбор методов селекции, кроссовера, мутации и стратегии замены популяции, а также критериев остановки, является важным фактором успеха.

\section{Представление решений (кодирование)}
\label{sec13} 

Представление решений, или кодирование, определяет, как каждая возможная <<гипотеза>> или решение задачи будет представлена в виде хромосомы, с которой работает генетический алгоритм. Выбор метода кодирования критически важен, поскольку он влияет на эффективность работы генетических операторов и общую производительность алгоритма. В пособии рассмат\-риваются несколько основных методов кодирования:

\begin{enumerate}
	\item \textit{бинарное кодирование}, в котором каждая хромосома представляется в виде бинарной строки (последовательности нулей и единиц). Данный метод кодирования широко используется благодаря своей простоте и универсальности. Однако, оно может быть неэффективным для задач, где решения имеют сложную структуру или требуют высокой точности представления;
	\item \textit{вещественное кодирование}, где каждая хромосома представляется в виде строки вещественных чисел. Данный метод кодирования хорошо подходит для задач, где решения представляют собой значения непрерывных переменных. Он позволяет более точно представлять решения и упрощает реализацию некоторых генетических операторов;
	\item \textit{перестановочное кодирование}, в котором хромосома представ\-ляется в виде перестановки элементов. И данный метод кодирования использует\-ся для задач, где порядок элементов имеет значение, например, для задачи коммивояжера;
	\item и другие методы кодирования.
\end{enumerate}

В частности, к задаче коррекции баз данных можно использовать следующие подходы к кодированию \texttt{SQL}-запросов: строковое и древовидное предс\-тавления. В первом случае \texttt{SQL}-запрос можно представить в виде строки символов. Этот метод прост в реализации, но требует специальных операторов, учитывающих синтаксис \texttt{SQL}. А во втором -- представить \texttt{SQL}-запрос в виде дерева, где узлы -- это операторы, функции, атрибуты и значения. Данный способ позволяет более эффективно применять генетические операторы, сохраняя синтаксическую корректность запросов.

Выбор методов кодирования оказывает существенное влияние на эффективность ГА. Важно тщательно анализировать особенности задачи и выбирать методы, которые наилучшим образом подходят для представления решений и манипулирования ими. В задаче коррекции баз данных необходимо учитывать синтаксис \texttt{SQL} и обеспечивать создание синтаксически корректных запросов после применения генетических операторов.


\chapter{Реализация алгоритма коррекции БД}
\label{chap2}
\section{Общее описание системы}
\label{sec21}

Разработанная система предназначена для коррекции ошибок в базах данных, возникающих из-за нарушения ограничений целостности. Система пост\-роена на основе генетического алгоритма и состоит из нескольких ключевых компонентов, обеспечивающих его эффективную работу в контексте исправления ошибок в базах данных. 

Основными компонентами системы являются:
\begin{enumerate}
	\item \textit{генетический алгоритм};
	\item \textit{парсер базы данных и ограничений целостности};
	\item \textit{модуль анализа ошибок};
	\item \textit{модуль генерации {SQL}-запросов};
	\item \textit{фитнес-функция};
	\item \textit{пользовательский интерфейс}.
\end{enumerate}

\textit{Генетический алгоритм} -- является ядром системы и отвечает за поиск оптимальных решений для исправления ошибок. Он реализует основные этапы генетического поиска: инициализация популяции, оценку фитнеса, селекцию, кроссовер, мутацию и замену популяции.

\textit{Парсер} -- отвечает за анализ структуры базы данных и извлечение информации об ограничениях целостности (например, \texttt{UNIQUE}, \texttt{NOT NULL}, \texttt{FOREIGN KEY}). Эта информация используется для генерации начальной популяции \texttt{SQL}-запросов и оценки их соответствия схеме базы данных.

\textit{Модуль анализа ошибок} -- анализирует сообщения об ошибках, генерируемые СУБД при выполнении \texttt{SQL}-запросов. Этот модуль определяет тип ошибки и местоположение ошибки (название таблицы, колонки, значение).

\textit{Модуль генерации {SQL}-запросов} -- отвечает за генерацию и модификацию {SQL}-запросов, используемых в генетическом алгоритме. Этот модуль генерирует как начальную популяцию запросов, так и новые запросы в процессе работы генетических операторов.

\textit{Фитнес-функция} -- оценивает качество {SQL}-запросов с точки зрения их соответствия требованиям целостности данных и эффективности исправления ошибок. Фитнес-функция учитывает несколько критериев: соблюдение ограничений целостности, соответствие схеме базы данных, эффективность исправления и сохранение данных.

И наконец, \textit{пользовательский интерфейс} -- предоставляет пользователю возможность взаимодействовать с системой, настраивать параметры генетического алгоритма, просматривать результаты работы и выбирать оптимальное решение для исправления ошибки.

Работа системы начинается с анализа базы данных и извлечения информации об ограничениях целостности с помощью парсера БД. Затем случайным образом генерируется начальная популяция \texttt{SQL}-запросов, с учетом информации о схеме БД и ограничений целостности. В процессе работы ГА \texttt{SQL}-запросы анализируются на наличие ошибок, оцениваются фитнес-функцией, отбираются лучшие решения, и генерируются новые запросы с помощью генетических операторов. Этот процесс повторяется до достижения критериев остановки.

\section{Фитнес-функция}
\label{sec22}

Оценка качества \texttt{SQL}-запросов осуществляется с помощью фитнес-функции, которая возвращает значение в диапазоне от 0 до 1. Значение 1 соот\-ветствует наилучшему решению, полностью удовлетворяющему всем критериям, а значение 0 – наихудшему. Фитнес-функция учитывает нормализированные критерии, представленные ниже.
\begin{enumerate}
	\item \textit{Соблюдение ограничений целостности}.
	
	Оценивается, не нарушает ли \texttt{SQL}-запрос существующие ограничения целостности в базе данных. Если нарушает ограничение, то назна\-чается штраф, уменьшающий итоговое значение. Критерий оценивается сле\-дующим образом:
	$$F_{integrity}(q)=1-\left(\frac{N_{violations}}{N_{maxViolations}}\right),$$
	где $N_{violations}$ -- количество нарушенных ограничений;
	
	$N_{maxViolations}$ -- максимальное возможное количество нарушений.
	
	\item \textit{Соответствие схеме базы данных}.
	
	Оценивается, соответствуют ли имена столбцов и типы данных, используемые в \texttt{SQL}-запросе, схеме БД. Критерий оценивается следующим образом:
	$$F_{schema}(q)=1-\left(\frac{N_{mismatches}}{N_{maxMismatches}}\right),$$
	где $N_{mismatches}$ -- количество несоответствий схеме;
	
	$N_{maxMismatches}$ -- максимальное возможное количество несоответствий.
	
	\item \textit{Эффективность исправления}.
	
	Оценивается, насколько эффективно \texttt{SQL}-запрос исправляет ошибку, минимизируя при этом изменения в базе данных. Критерий оценивает\-ся на основе количества измененных строк:
	$$F_{efficiency}(q)=\left(\frac{1}{1+ln(1+N_{modified})}\right),$$
	где $N_{modified}$ -- количество строк, измененных запросом.
	
	Эта формула обеспечивает, что чем меньше строк изменено, тем ближе значение к 1. Логарифм используется для сглаживания влияния больших значений $N_{modified}$.
	
	\item \textit{Сохранение данных}.
	
	Оценивается, насколько \texttt{SQL}-запрос сохраняет данные в БД. Запросы, приводящие к удалению данных, получают более низкую оценку. Этот критерий можно реализовать на основе типа операции:
	$$F_{efficiency}(q)=OperationWeight,$$
	где $OperationWeight$ -- вес операции, принимающий значения в диапазоне от 0 до 1 в зависимости от типа операции:
	\begin{itemize}
		\item \texttt{UPDATE} и \texttt{INSERT}: 1 (наилучшее, обновляет или сохраняет данные);
		\item \texttt{DELETE}: 0.25 (наихудшее, удаляет данные).
	\end{itemize}
	
	\item \textit{Использование заглушек}.
	
	Оценивается временных значений для обхода ограничений, например, при нарушении \texttt{UNIQUE}. Использование заг\-лушек не является идеальным решением, поэтому за это назначается штраф. Критерий оценивается следующим образом:
	$$F_{placeholder}(q)=1-\left(\frac{N_{placeholder}}{N_{maxPlaceholder}}\right),$$
	где $N_{placeholder}$ -- количество использованных заглушек;
	
	$N_{maxPlaceholder}$ -- максимальное возможное количество заглушек (определяется исходя из количества полей, требующих уникальных значений).
\end{enumerate}

Фитнес-функция представляет собой взвешенную сумму этих критериев:
\begin{equation}
	\begin{split}
		F(q) = {} & w_1 F_{integrity}(q) + w_2 F_{schema}(q) + w_3 F_{efficiency}(q) + \\
		& w_4 F_{preservation}(q) + w_5 F_{placeholder}(q),
	\end{split}
\end{equation}
где $F(q)$ -- значение фитнес-функции для \texttt{SQL}-запроса $q$ (от 0 до 1);

$F_{integrity}$ -- оценка соблюдения ограничений целостности для запроса $q$;

$F_{schema}$ -- оценка соответствия схеме базы данных для запроса $q$;

$F_{efficiency}$ -- оценка эффективности исправления ошибки для запроса $q$;

$F_{prevervation}$ -- оценка сохранения данных для запроса $q$;

$F_{placeholder}$ -- оценка использования заглушек для запроса $q$;

$w_1,w_2,w_3,w_4,w_5$ -- веса, определяющие важность каждого критерия (сумма весов должна быть равна 1).

\section{Выбор и реализация генетических операторов и параметров}
\label{sec23}

При разработке генетического алгоритма для коррекции ошибок в базах данных был сделан ряд конкретных выборов относительно используемых генетических операторов, параметров алгоритма и других деталей реализации. Эти выборы были сделаны на основе анализа предметной области, результатов предварительных экспериментов и соображений вычислительной эффективности.

\texttt{SQL}-запросы в разработанной системе представляются в виде строк. Этот выбор был сделан в первую очередь для простоты реализации и гибкости. Строковое представление позволяет легко представлять запросы различной сложности, включая \texttt{INSERT}, \texttt{UPDATE} и \texttt{DELETE}. Такой подход обеспечивает возможность создания разнообразных генетических мутаций и кроссоверов, которые будут рассмотрены далее.

Строковое представление (\texttt{SQL}-запрос) обладает следующей структурой:

\texttt{UPDATE customers SET email = 'new.email@example.com'}

\texttt{WHERE customer\_id = 1;}

В строковом представлении он будет представлен полностью.

Для обеспечения синтаксической корректности генерируемых \texttt{SQL}-выраже\-ний используется специализированный парсер и генератор, который позво\-ляет описать грамматику \texttt{SQL} в формальном виде.

Парсер выполняет следующие проверки:
\begin{enumerate}
\item проверка синтаксической корректности ключевых слов и символов (к примеру, правильность написания \texttt{UPDATE, SET, WHERE} и т.д.);
\item проверка наличия и корректности имен таблиц и столбцов (например, проверка существования таблицы \texttt{customers} и столбца \texttt{email}). Для этого парсер взаимодействует с метаданными базы данных;
\item проверка соответствия типов данных при присваивании значений (например, попытка присвоения текстового значения полю типа \texttt{INTEGER} будет обнаружена);
\item проверка соблюдения ограничений \texttt{UNIQUE, NOT NULL} (на уровне синтаксического анализа, а не выполнения запроса);
\item обработка выражений в \texttt{WHERE} предложении.
\end{enumerate}
В случае обнаружения синтаксической ошибки парсер генерирует сообщение об ошибке, которое используется для отбрасывания некорректного потомка и генерации нового.

Альтернативным подходом является представление \texttt{SQL}-запросов в виде абстрактных синтаксических деревьев (\texttt{AST}) или графов. Преимуществом такого подхода является более удобная работа с элементами запроса и упрощенный процесс генерации. Однако, это усложняет реализацию операторов генетического алгоритма, а также требует более сложного процесса парсинга и генерации \texttt{SQL}-кода. Было принято решение о использовании строкового представления запросов, так как это обеспечило простоту реализации и возможность быстрого прототипирования.

Для отбора наиболее приспособленных особей из популяции использовалась турнирная селекция с размером турнира \texttt{k = 5}. Данный вид селекции был выбран из-за ее простоты реализации, высокой скорости работы и способности поддерживать разнообразие популяции по сравнению с другими методами селекции. При турнирной селекции случайным образом выбирается \texttt{k} особей из популяции, и из них особь с наилучшим значением фитнес-функции переходит на следующий этап, а остальные отсеиваются. Этот процесс повторяется для формирования популяции следующего поколения.

Выбор размера турнира \texttt{k = 5} был сделан на основе предварительных экспериментов. Небольшие значения \texttt{k} (например, 2 или 3) приводили к недос\-таточному селективному давлению, что замедляло сходимость алгоритма. Большие значения \texttt{k} (например, 10 или 15) приводили к слишком сильному селективному давлению, что снижало разнообразие популяции и увеличивало вероятность преждевременной сходимости. Было проведено несколько серий экспериментов с различными значениями \texttt{k} (2, 3, 5, 7, 10). На основе полученных результатов было установлено, что значение \texttt{k = 5} обеспечивает оптимальный баланс между скоростью сходимости и сохранением разнообразия популяции.

Альтернативными методами селекции были: пропорциональная селекция (рулетка) и ранговая селекция. Пропорциональная была отклонена из-за ее чувствительности к выбросам в значениях фитнес-функции, а ранговая -- из-за меньшей эффективности по сравнению с турнирной селекцией.

В качестве оператора кроссовера использовался одноточечный кроссовер. Это один из наиболее распространенных операторов кроссовера, который зак\-лючается в случайном выборе точки разрыва в строковом представлении \texttt{SQL}-запроса и обмене частями запросов между двумя родителями.

Процесс одноточечного кроссовера можно описать следующим образом:
\begin{enumerate}
\item выбираются два родителя (особи с лучшими значениями фитнес-функ\-ции) для кроссовера;
\item случайным образом определяется точка разрыва в строковом представлении \texttt{SQL}-запроса;
\item осуществляется обмен частями \texttt{SQL}-запросов между родителями, начиная с точки разрыва.
\end{enumerate}
Например:

родитель 1: \texttt{UPDATE products SET price = 100 WHERE product\_id = 1;}

родитель 2: \texttt{UPDATE products SET price = 55 WHERE product\_id = 6;}

Точка разрыва после \texttt{WHERE} для первого родителя.

потомок 1: \texttt{UPDATE products SET price = 100 WHERE product\_id = 6;}

потомок 2: \texttt{UPDATE products SET price = 55 WHERE product\_id = 1;}

Перед применением оператора кроссовера выполняется проверка синтаксической корректности получающихся потомков с помощью ранее описанного парсера.
Если потомок оказывается синтаксически некорректным, он отбрасывается и генерируется заново. Вероятность генерации корректного потомка составляет 80\%.
Этот выбор был сделан из-за его простоты реализации и эффективности. Одноточечный кроссовер позволяет создавать новые запросы путем комбинирования частей существующих, что обеспечивает разнообразие популяции и позволяет алгоритму исследовать различные варианты решений.

Рассматривались также двухточечный и равномерный кроссовер. Двухточечный предполагает наличие двух точек разрыва, что потенциально может привести к более разнообразным результатам. Однако, это увеличивает ве\-роятность генерации синтаксически некорректных запросов, что требует более сложной обработки ошибок, а равномерный кроссовер предполагает обмен отдельными атрибутами между родителями, что может быть более эффективно для решения специфических задач, но, в нашем случае, может привести к потере взаимосвязей между атрибутами запроса.

Для внесения случайных изменений в \texttt{SQL}-запросы использовался оператор мутации. Он выбирает случайным образом атрибут (например, имя таблицы, имя столбца, значение) и заменяет его на другое допустимое значение. Процесс мутации включает следующие этапы:
\begin{enumerate}
\item случайным образом выбирается особь для мутации;
\item определяется атрибут, который будет подвергнут мутации;
\item генерируется новое значение для выбранного атрибута.
\end{enumerate}
Примеры мутаций:

изменение имени таблицы: \texttt{UPDATE customers ... -> UPDATE users ...};

изменение имени столбца: \texttt{SET email = ... -> SET phone = ...};

изменение значения: \texttt{WHERE customer\_id = 1 -> WHERE customer\_id = 2};

Примеры мутаций для разных типов данных (числа, строки, даты):
\begin{itemize}
\item \texttt{total\_amount = 100.00 -> total\_amount = 120.00};
\item \texttt{email = 'test@example.com' -> email = 'new@example.com'};
\item \texttt{order\_date = '2024-05-15' -> order\_date = '2024-06-01'}.
\end{itemize}

После каждой мутации выполняется проверка синтаксической корректнос\-ти нового \texttt{SQL}-запроса с помощью парсера. Если запрос некорректен, мутация отменяется и генерируется новая. Также проверяется соблюдение ограничений целостности.

Процесс генерации нового значения зависит от типа мутируемого атрибута. Для атрибутов, имеющих ограниченное число допустимых значений, новое значение выбирается случайным образом из списка допустимых значений; для числовых атрибутов -- новое значение генерируется в определенном диапазоне; для строковых -- новое значение может быть сгенерировано случайным образом или получено из существующих данных.

Рассматривались альтернативные методы мутации, такие как мутация, основанная на операциях над деревьями (AST), но они потребовали бы более сложной реализации и не обеспечили бы значительного преимущества по сравнению с выбранным подходом. Также рассматривалась мутация, в которой мутировались бы части запроса, но это привело бы к низкой вероятности генерации корректного запроса.

Не существует универсального решения для распределения весов в фитнес-функции, поэтому был выбран способ с ориентиром на важность целостности системы и весов:
\begin{itemize}
\item $w_1 (F_{integrity})= 0.4$;
\item $w_2 (F_{schema})= 0.3$;
\item $w_3 (F_{efficiency})= 0.1$;
\item $w_4 (F_{preservation})= 0.1$;
\item $w_5 (F_{placeholder)}= 0.1$.
\end{itemize}

И так же после серии экспериментов были выбраны следующие параметры алгоритма:
\begin{itemize}
	\item размер популяции -- 50 особей;
	\item вероятность кроссовера -- 0.8;
	\item вероятность мутации -- 0.1;
	\item количество лучших решений -- 1;
	\item критерий остановки -- достижение максимального количества поколений (100).
\end{itemize}

\section{Интерфейс программы}
\label{sec24}

\begin{figure}
	\centering
%	\includegraphics[width=0.9\textwidth]{interface.jpg}
	\caption{Интерфейс программы.}
	\label{fig1}
\end{figure}
Разработанная программа имеет графический пользовательский интерфейс (\texttt{GUI}), который позволяет пользователю взаимодействовать с системой и анализировать результаты ее работы.
Интерфейс программы состоит из следующих основных элементов:
\begin{enumerate}
\item \textit{панель авторизации БД}. Позволяет пользователю ввести параметры подключения к базе данных (\texttt{URL}, имя пользователя, пароль);
\item \textit{поле ввода {SQL}-запроса}. Позволяет пользователю ввести \texttt{SQL}-запрос, который необходимо выполнить и исправить (если он приводит к ошибке);
\item \textit{панель настройки параметров генетического алгоритма}. Позволяет пользователю настроить параметры ГА (размер популяции, вероятность мутации и кроссовера, количество выводимых лучших решений и поколений);
\item \textit{кнопка «Execute»}. Запускает процесс выполнения генетического алгоритма для исправления ошибки;
\item \textit{панель отображения ошибок}. Отображает информацию о невозможности выполнить \texttt{SQL}-запрос и текст ошибки;
\item \textit{панель отображения результатов}. Демонстрирует результаты оптимизации:
\begin{itemize}
\item список предложенных \texttt{SQL}-запросов для исправления ошибки (в порядке убывания значения фитнес-функции);
\item количественную оценку пригодности каждого запроса через фитнес-функцию.
\end{itemize}
\end{enumerate}
Вид интерфейса программы представлен на рис. \ref{fig1}.


\chapter{Экспериментальные результаты}
\label{chap3}
\section{Тестовые данные}
\label{sec31}

\begin{figure}
	\centering
%	\includegraphics[width=0.9\textwidth]{bd.jpg}
	\caption{Тестовая база данных.}
	\label{fig2}
\end{figure}

Для экспериментальной оценки разработанного генетического алгоритма была создана тестовая база данных, моделирующая структуру и ограничения реальных баз данных. Схема БД представлена на рис. \ref{fig2}.

Таблицы содержат различные типы данных (\texttt{INTEGER}, \texttt{VARCHAR}, \texttt{DECIMAL}, \texttt{DATE}) и ограничения целостности: \texttt{PRIMARY KEY}, \texttt{FOREIGN KEY}, \texttt{UNIQUE}, \texttt{NOT NULL}, \texttt{CHECK}.

Сама же база данных моделирует структуру интернет-магазина и состоит из следующих таблиц:
\begin{enumerate}
\item \texttt{customers} -- хранит информацию о зарегистрированных покупателях магазина.

Атрибуты:
\begin{itemize}
\item \texttt{customer\_id} (\texttt{SERIAL PRIMARY KEY}): автоматически генерируемый уникальный идентификатор покупателя (первичный ключ);
\item \texttt{first\_name} (\texttt{VARCHAR(50) NOT NULL}): имя покупателя;
\item \texttt{last\_name} (\texttt{VARCHAR(50) NOT NULL}): фамилия покупателя;
\item \texttt{email} (\texttt{VARCHAR(100) UNIQUE NOT NULL}): адрес электронный поч\-ты покупателя (уникальный);
\item \texttt{phone} (\texttt{VARCHAR(20)}): номер телефона покупателя;
\item \texttt{address} (\texttt{VARCHAR(200)}): адрес покупателя;
\item \texttt{city} (\texttt{VARCHAR(50)}): город покупателя;
\item \texttt{state} (\texttt{VARCHAR(50)}): область покупателя;
\item \texttt{zip\_code} (\texttt{VARCHAR(10)}): почтовый индекс покупателя;
\item \texttt{create\_at} (\texttt{TIMESTAMP DEFAULT NOW()}): дата и время регистрации покупателя.
\end{itemize}

\item \texttt{orders} -- хранит информацию о заказах, сделанных покупателями.

Атрибуты:
\begin{itemize}
	\item \texttt{order\_id} (\texttt{SERIAL PRIMARY KEY}): автоматически генерируемый уникальный идентификатор заказа (первичный ключ);
	\item \texttt{customer\_id} (\texttt{INTEGER NOT NULL}): идентификатор покупателя, сделавшего заказ (является внешним ключом, ссылающимся на таблицу \texttt{customers}, обеспечивая связь между заказами и покупателями);
	\item \texttt{order\_date} (\texttt{TIMESTAMP DEFAULT NOW()}): дата и время оформления заказа;
	\item \texttt{shipping\_address} (\texttt{VARCHAR(200)}): адрес доставки заказа;
	\item \texttt{total\_amount} (\texttt{DECIMAL(10, 2)} \texttt{CHECK (total\_amount > 0)} \texttt{NOT} \\\texttt{NULL}): общая сумма заказа (должна быть больше нуля);
	\item \texttt{status} (\texttt{VARCHAR(20) CHECK (status IN (‘PROCESSING’, ‘DELI\-VERED’, PENDING’, ‘CANCELLED’, ‘SHIPPED’))}): статус заказа, а именно -- <<Обработка>>, <<Доставлен>>, <<Ожидание>>, <<Отправлен>>, <<Отменен>> (\texttt{CHECK}-ограничение гарантирует, что статус заказа будет одним из допустимых значений);
	\item \texttt{shipping\_method\_id} (\texttt{INTEGER}): идентификатор метода доставки, использованного для заказа (является внешним ключом, ссылающимся на таблицу \texttt{shipping\_methods}).
\end{itemize}

\item \texttt{shipping\_methods} -- хранит информацию о доступных методах доставки.

Атрибуты:
\begin{itemize}
	\item \texttt{shipping\_method\_id} (\texttt{SERIAL PRIMARY KEY}): автоматически генерируемый уникальный идентификатор метода доставки (первичный ключ);
	\item \texttt{name} (\texttt{VARCHAR(50) NOT NULL}): название метода доставки;
	\item \texttt{cost} (\texttt{DECIMAL(5, 2) NOT NULL}): стоимость доставки.
\end{itemize}

\item \texttt{products} -- хранит информацию о продуктах, предлагаемых в магазине.

Атрибуты:
\begin{itemize}
	\item \texttt{product\_id} (\texttt{SERIAL PRIMARY KEY}): автоматически генерируемый уникальный идентификатор продукта (первичный ключ);
	\item \texttt{name} (\texttt{VARCHAR(100) NOT NULL}): название продукта;
	\item \texttt{description} (\texttt{TEXT}): описание товара;
	\item \texttt{price} (\texttt{DECIMAL(10, 2) NOT NULL CHECK (price > 0)}): цена продукта (\texttt{CHECK}-ограничение гарантирует что цена будет больше нуля);
	\item \texttt{categoty\_id} (\texttt{INTEGER NOT NULL}): идентификатор категории продукта (является внешним ключом, ссылающимся на таблицу \texttt{cate\-gories}).
\end{itemize}

\item \texttt{categories} -- хранит информацию о категориях продуктов.

Атрибуты:
\begin{itemize}
	\item \texttt{category\_id} (\texttt{SERIAL PRIMARY KEY}): автоматически генерируемый уникальный идентификатор категории (первичный ключ);
	\item \texttt{name} (\texttt{VARCHAR(50) NOT NULL UNIQUE}): название категории (уникальное).
\end{itemize}

\item \texttt{order\_items} -- хранит информацию о товарах, входящих в конкретный заказ (связывает таблицу заказов и продуктов).

Атрибуты:
\begin{itemize}
	\item \texttt{order\_item\_id} (\texttt{SERIAL PRIMARY KEY}): генерируемый автоматичес\-ки уникальный идентификатор элемента заказа (первичный ключ);
	\item \texttt{order\_id} (\texttt{INTEGER NOT NULL}): идентификатор заказа, к которому относится элемент (является внешним ключом, ссылающимся на таблицу \texttt{orders});
	\item \texttt{product\_id} (\texttt{INTEGER NOT NULL}): идентификатор продукта, входящему в заказ (является внешним ключом, ссылающимся на таблицу \texttt{products});
	\item \texttt{quanty} (\texttt{INTEGER NOT NULL CHECK (quantity > 0)}): кол-во данного продукта в заказе (\texttt{CHECK}-ограничение гарантирует что количество будет больше нуля);
	\item \texttt{price} (\texttt{DECIMAL(10, 2) NOT NULL CHECK (price > 0)}): стоимость единицы товара на момент заказа (\texttt{CHECK}-ограничение гарантирует что цена будет больше нуля).
\end{itemize}

\item \texttt{product\_reviews} -- хранит отзывы покупателей о продуктах.

Атрибуты:
\begin{itemize}
	\item \texttt{review\_id} (\texttt{SERIAL PRIMARY KEY}): уникальный автоматически генерируемый идентификатор отзыва (первичный ключ);
	\item \texttt{product\_id} (\texttt{INTEGER NOT NULL}): идентификатор продукта, о котором написан отзыв (является внешним ключом, ссылающимся на таблицу \texttt{products});
	\item \texttt{customer\_id} (\texttt{INTEGER NOT NULL}): идентификатор покупателя, написавшего отзыв (является внешним ключом, ссылающимся на таблицу \texttt{customers});
	\item \texttt{rating} (\texttt{INTEGER CHECK (rating >= 1 AND rating <= 5)}): оценка продукта от 1 до 5 (\texttt{CHECK}-ограничение гарантирует что оценка будет в допустимом диапазоне);
	\item \texttt{comment} (\texttt{TEXT}): текст отзыва;
	\item \texttt{create\_at} (\texttt{TIMESTAMP DEFAULT NOW()}): дата и время создания отзыва.
\end{itemize}

\item \texttt{shopping\_cart} -- хранит информацию о корзинах покупок покупателей.

Атрибуты:
\begin{itemize}
	\item \texttt{cart\_id} (\texttt{SERIAL PRIMARY KEY}): автоматически генерируемый уникальный идентификатор корзины (первичный ключ);
	\item \texttt{customer\_id} (\texttt{INTEGER NOT NULL}): идентификатор покупателя, которому принадлежит корзина (является внешним ключом, ссылаю\-щимся на таблицу \texttt{customers});
	\item \texttt{create\_at} (\texttt{TIMESTAMP DEFAULT NOW()}): дата и время создания корзины.
\end{itemize}

\item \texttt{cart\_items} -- хранит информацию о товарах, находящихся в корзине покупок (связующая таблица между корзинами и продуктами).

Атрибуты:
\begin{itemize}
	\item \texttt{cart\_item\_id} (\texttt{SERIAL PRIMARY KEY}): уникальный авто-генерируе\-мый идентификатор элемента корзины (первичный ключ);
	\item \texttt{cart\_id} (\texttt{INTEGER NOT NULL}): идентификатор корзины, к которому относится элемент (является внешним ключом, ссылающимся на таблицу \texttt{shopping\_cart});
	\item \texttt{product\_id} (\texttt{INTEGER NOT NULL}): идентификатор продукта, находящегося в корзине (является внешним ключом, ссылающимся на таблицу \texttt{products});
	\item \texttt{quanty} (\texttt{INTEGER NOT NULL CHECK (quantity > 0)}): кол-во данного продукта в корзине (\texttt{CHECK}-ограничение гарантирует что количество будет больше нуля).
\end{itemize}
\end{enumerate}

\section{Проведение экспериментов и их анализ}
\label{sec32}

Для оценки эффективности разработанного генетического алгоритма была проведена серия экспериментов с использованием тестовой базы данных и различных сценариев генерации ошибок.

\begin{figure}
	\centering
%	\includegraphics[width=0.9\textwidth]{error.jpg}
	\caption{Окно об ошибке из примера 1.}
	\label{fig3}
\end{figure}

\textit{Пример 1.} Рассмотрим \texttt{SQL}-запрос на обновление данных:

\texttt{UPDATE customers} 

\texttt{SET email =’ivan.ivanov@example.com’}

\texttt{WHERE id = 2;}

После запуска алгоритма появляется окно об ошибке, представленное на рис. \ref{fig3}. Ошибка гласит, что такое значение email уже существует в таблице, а это ведёт к нарушению ограничения целостности \texttt{UNIQUE}. 

В таблице \ref{table1} описан результат работы программы.

\begin{table}
	\centering
	\begin{tabularx}{\linewidth}{@{}>{\hsize=0.5\hsize}X>{\hsize=1.5\hsize}X@{}}
		\toprule
		\textbf{Решение} & \textbf{SQL-запросы} \\
		\midrule
		Пример 1 
		
		(Фитнес: 0.86) & 
		Нарушено \texttt{UNIQUE} ограничение для поля \texttt{email}.
		
		Вставлена заглушка: " Дубликат данных".\\
		& Итоговый \texttt{SQL}-запрос: \\
		& \texttt{UPDATE customers SET email = ‘Дубликат данных' WHERE id = 2;} \\
		\bottomrule
	\end{tabularx}
	\caption{Результаты работы программы из примера 1.}
	\label{table1}
\end{table}

Как видим, результат фитнеса высок, программа подставила заглушку, так как хромосомы создаются на основе данных в таблице, а для текстового значения создавать случайную запись не совсем корректно, поэтому программа указывает, что данные по \texttt{email} нуждаются в замене.

\textit{Пример 2.} Рассмотрим \texttt{SQL}-запрос на добавление данных:

\texttt{INSERT INTO customers (first\_name, last\_name)} 

\texttt{VALUES (‘Новый’, ’Клиент’);}

После запуска алгоритма появляется окно об ошибке, которое информирует о том, что поле \texttt{email} не может быть пустым, а именно срабатывает ограничение \texttt{NOT NULL}. В таблице \ref{table2} описан результат работы программы.

\begin{table}
	\centering
	\begin{tabularx}{\linewidth}{@{}>{\hsize=0.5\hsize}X>{\hsize=1.5\hsize}X@{}}
		\toprule
		\textbf{Решение} & \textbf{SQL-запросы} \\
		\midrule
		Пример 2 
		
		(Фитнес: 0.86) & 
		Нарушено \texttt{NOT NULL } ограничение для поля \texttt{email}.
		
		Вставлена заглушка: " Нехватка данных".\\
		& Итоговый \texttt{SQL}-запрос: \\
		& \texttt{INSERT INTO customers (first\_name, last\_name, email) VALUES (‘Новый’, ’Клиент’, ‘Нехватка данных’);} \\
		\bottomrule
	\end{tabularx}
	\caption{Результаты работы программы из примера 2.}
	\label{table2}
\end{table}


В данном случае произошла аналогичная ситуация, алгоритм так же воспользовался заглушкой, чтобы не генерировать случайную комбинацию символов в поле, поэтому и результаты схожи.

\textit{Пример 3.} Рассмотрим \texttt{SQL}-запрос на удаление данных:

\texttt{DELETE FROM customers WHERE id = 1;}

После запуска алгоритма появляется окно об ошибке, которое информи\-рует о том, что нарушено ограничение \texttt{FOREIGN KEY}, так как существуют записи в других таблицах, использующие информацию об покупателе с идентификатором 1. В таблице \ref{table3} представлен результат работы программы.
	
\begin{table}
	\centering
	\begin{tabularx}{\linewidth}{@{}>{\hsize=0.5\hsize}X>{\hsize=1.5\hsize}X@{}}
		\toprule
		\textbf{Решение} & \textbf{SQL-запросы} \\
		\midrule
		Пример 3
		
		Решение 1 
		
		(Фитнес: 0.76) & 
		Обнаружены связанные записи в \texttt{orders} (FOREIGN KEY: \texttt{customer\_id}):
		
		Итоговые \texttt{SQL}-запросы: \\
		& \texttt{UPDATE orders SET customer\_id = 2 WHERE customer\_id = 1;} \\
		& \texttt{UPDATE product\_reviews SET customer\_id = 2 WHERE customer\_id = 1;} \\
		& \texttt{DELETE FROM customers WHERE id = 1;} \\
		\cmidrule{1-2}
		Пример 3
		
		Решение 2 
		
		(Фитнес: 0.40) & 
		Обнаружены связанные записи в \texttt{orders} (FOREIGN KEY: \texttt{customer\_id}):
		
		Итоговые \texttt{SQL}-запросы: \\
		& \texttt{DELETE FROM orders WHERE customer\_id = 1;} \\
		& \texttt{DELETE FROM customers WHERE id = 1;} \\
		\bottomrule
	\end{tabularx}
	\caption{Результаты работы программы из примера 3.}
	\label{table3}
\end{table}

\begin{table}
	\centering
	\begin{tabularx}{\linewidth}{@{}>{\hsize=0.5\hsize}X>{\hsize=1.5\hsize}X@{}}
		\toprule
		\textbf{Решение} & \textbf{SQL-запросы} \\
		\midrule
		Пример 4 
		
		(Фитнес: 0.96) & 
		Нарушено \texttt{CHECK} ограничение для поля \texttt{rating}.
		
		Итоговый \texttt{SQL}-запрос: \\
		& \texttt{INSERT INTO product\_reviews (customer\_id, product\_id, rating, review\_date, comment) VALUES (1, 1, 5, '2024-10-27', 'Отличный продукт!');} \\
		\bottomrule
	\end{tabularx}
	\caption{Результаты работы программы из примера 4.}
	\label{table4}
\end{table}

\begin{table}
	\centering
	\begin{tabularx}{\linewidth}{@{}>{\hsize=0.5\hsize}X>{\hsize=1.5\hsize}X@{}}
		\toprule
		\textbf{Решение} & \textbf{SQL-запросы} \\
		\midrule
		Пример 5 
		
		(Фитнес: 0.96) & 
		Нарушен тип данных для поля \texttt{price}.
		
		Итоговый \texttt{SQL}-запрос: \\
		& \texttt{UPDATE products SET price = 153.26 WHERE product\_id = 1;} \\
		\bottomrule
	\end{tabularx}
	\caption{Результаты работы программы из примера 5.}
	\label{table5}
\end{table}

\begin{table}
	\centering
	\begin{tabularx}{\linewidth}{@{}>{\hsize=0.5\hsize}X>{\hsize=1.5\hsize}X@{}}
		\toprule
		\textbf{Решение} & \textbf{SQL-запросы} \\
		\midrule
		Пример 6
		
		Решение 1
		
		(Фитнес: 0.96) & 
		Нарушено \texttt{FOREIGN KEY} ограничение для поля \texttt{customer\_id}.
		
		Итоговые \texttt{SQL}-запросы: \\
		& \texttt{UPDATE orders SET customer\_id = 3, total\_amount = 585.00 WHERE order\_id = 14;} \\
		\cmidrule{1-2}
		Пример 6
		
		Решение 2
		
		(Фитнес: 0.96) & 
		Нарушено \texttt{FOREIGN KEY} ограничение для поля \texttt{customer\_id}.
		
		Итоговые \texttt{SQL}-запросы: \\
		& \texttt{UPDATE orders SET customer\_id = 5, total\_amount = 360.00 WHERE order\_id = 14;} \\
		\cmidrule{1-2}
		Пример 6
		
		Решение 3
		
		(Фитнес: 0.96) & 
		Нарушено \texttt{FOREIGN KEY} ограничение для поля \texttt{customer\_id}.
		
		Итоговые \texttt{SQL}-запросы: \\
		& \texttt{UPDATE orders SET customer\_id = 6, total\_amount = 887.00 WHERE order\_id = 14;} \\
		\bottomrule
	\end{tabularx}
	\caption{Результаты работы программы из примера 6.}
	\label{table6}
\end{table}

\begin{table}
	\centering
	\begin{tabularx}{\linewidth}{@{}>{\hsize=0.5\hsize}X>{\hsize=1.5\hsize}X@{}}
		\toprule
		\textbf{Решение} & \textbf{SQL-запросы} \\
		\midrule
		Пример 7
		
		(Фитнес: 0.92) & 
		Обнаружены связанные записи в \texttt{products} (FOREIGN KEY: \texttt{category\_id}):
		
		Итоговые \texttt{SQL}-запросы: \\
		& \texttt{UPDATE products SET category\_id = 3 WHERE category\_id = 4;} \\
		& \texttt{DELETE FROM categories WHERE category\_id = 4;} \\
		\bottomrule
	\end{tabularx}
	\caption{Результаты работы программы из примера 7.}
	\label{table7}
\end{table}

В первом результате программа предложила обновить идентификаторы покупателя, используемые в двух других таблицах и после этого удалить самого покупателя. Оценка фитнеса высокая, значит это решение валидно.

Если изменить уменьшить количество поколений до 7, а размер популяции до 3 особей, то получились результаты под вторым решением в таблице \ref{table3}.

В данном случае заметно, как сильно упала оценка фитнеса, что говорит о некорректности данного решения, что видно и по конечным запросам. В итоге вместо замены идентификаторов, алгоритм предложил удалить записи в таблице \texttt{orders}, где идентификатор покупателя равен 1, но тем самым не учитывается, что в этих записях могут быть связи с другими таблицами, а также совсем не учтена таблица \texttt{product\_reviews} из прошлого варианта. Следовательно, фитнес-функция правильно указала на невалидность данного решения.

\textit{Пример 4.} Рассмотрим \texttt{SQL}-запрос на добавление данных:

\texttt{INSERT INTO product\_reviews (customer\_id, product\_id,} 

\texttt{rating, review\_date, comment)}

\texttt{VALUES (1, 1, 6, '2024-10-27', 'Отличный продукт!');}

Этот запрос пытается добавить отзыв с рейтингом 6, что нарушает \texttt{CHECK}-ограничение, так как поле \texttt{rating} содержит в себе числа от 1 до 5. Результат работы программы представлен в таблице \ref{table4}.

В этом примере алгоритм предлагает изменить оценку рейтинга с 6 на 5, что будет удовлетворять ограничению \texttt{CHECK}.


\textit{Пример 5.} Рассмотрим \texttt{SQL}-запрос на обновление данных:

\texttt{UPDATE products SET price = 'abc' WHERE product\_id = 1;} 

В данном запросе нарушается тип данных, так как \texttt{price} -- числовое поле, а в запросе используется текстовое значение. Результат работы программы представлен в таблице \ref{table5}.


В этом примере текстовое значение было изменено на числовое, алгоритм находит его на основе имеющихся данных. И в итоге получилось устранить ошибку.

\textit{Пример 6.} Рассмотрим \texttt{SQL}-запрос на обновление данных:

\texttt{UPDATE orders SET customer\_id = 999, total\_amount = -50.00}

\texttt{WHERE order\_id = 14;}

В данном случае нарушается не только ограничение на внешний ключ к таблице \texttt{customers}, так как 999 пользователя не существует, но и \texttt{CHECK}-ограничение у поля \texttt{total\_amount}, значение которого должно быть больше нуля. В таблице \ref{table6} представлен результат работы программы.

В этом примере были выведены 3 лучших результата работы. Алгоритм заменил несуществующий идентификатор на тот, что есть в БД, поэтому ограничение внешнего ключа соблюдено, и изменил \texttt{total\_amount} на положительное число, используя другие имеющиеся в базе данных значениях данного поля, теперь ограничение \texttt{CHECK} не нарушается. Благодаря случайностям, которые выполняются в процессе эволюции, получились разнообразные результаты, все из которых могут использоваться для редактирования изначального невалидного \texttt{SQL}-запроса.

\textit{Пример 7.} Рассмотрим \texttt{SQL}-запрос на удаление данных:

\texttt{DELETE FROM categories WHERE category\_id = 4}

После запуска алгоритма появляется окно об ошибке, которое информи\-рует о том, что нарушено ограничение \texttt{FOREIGN KEY}, так как существует запись в таблице \texttt{products}, использующая информацию о категории товара с идентификатором 4. В таблице \ref{table7} представлен результат работы программы.

В данном примере наблюдается высокая оценка фитнеса, и она обоснована тем, что алгоритм предложил заменить идентификатор категории в таблице \texttt{products} на другой существующий, а после этого удалить категорию. В итоге ГА нашёл решение, которое соблюдает всем ограничениям целостности.

На основе представленных выше примеров можно провести следующий анализ эффективности разработанной фитнес-функции. 

При использовании оптимальных параметров генетического алгоритма (размер популяции: 50, вероятность кроссовера: 0.8, вероятность мутации: 0.1, количество поколений: 100), алгоритм смог успешно скорректировать все представленные в примерах \texttt{SQL}-запросы с ошибками. Это соответствует 100\% успешности коррекции для данного набора примеров.

Однако, следует отметить, что данный результат не позволяет гарантировать аналогичную успешность на более сложных и разнообразных наборах данных. При работе с реальными базами данных, содержащими более сложные запросы и ограничения, а также с более разнообразными типами ошибок, успешность коррекции может быть ниже.

Тем не менее, стоит отметить, что в ходе разработки и тестирования алгоритма были рассмотрены и другие примеры, не вошедшие в данное описание, и алгоритм успешно справлялся и с ними, что подтверждает его общую работоспособность и применимость к широкому кругу задач коррекции \texttt{SQL}-запросов.

Для подтверждения эффективности разработанной фитнес-функции была проведена серия экспериментов с использованием набора тестовых данных. В результате анализа значений фитнес-функции для корректных (среднее = 0.878) и некорректных (среднее = 0.449) \texttt{SQL}-запросов, был проведен \texttt{t}-тест, который показал статистически значимую разницу между группами \texttt{(t = 12.44, p < 0.001)}. Данный результат свидетельствует о способности фитнес-функции эффективно отличать валидные решения от невалидных, что является необходимым условием для успешной работы генетического алгоритма.

Так же был проведен дополнительный анализ для оценки эффективнос\-ти фитнес-функции в зависимости от типа \texttt{SQL}-запроса (\texttt{INSERT, UPDATE, DELETE}). Получены следующие результаты:
\begin{enumerate}
\item для операций \texttt{INSERT} и \texttt{UPDATE} средним значением фитнес-функции для корректных запросов является 0.95. Это указывает на высокую точность фитнес-функции при оценке данных операций, что позволяет генетическому алгоритму эффективно находить оптимальные решения для коррекции ошибок при вставке и обновлении данных;
	
\item для операций \texttt{DELETE} средним же значением фитнес-функции для корректных запросов является 0.79, что несколько более низкое значение по сравнению с \texttt{INSERT} и \texttt{UPDATE}. Это объясняется тем, что операции \texttt{DELETE} часто связаны с удалением связанных данных, что требует учета дополнительных факторов, таких как сохранение целостности базы данных и минимизация потери информации. Фитнес-функция в данном случае должна находить баланс между устранением ошибки и сохранением ценных данных, что может приводить к менее высоким значениям фитнес-функции, но при этом обеспечивать приемлемое качество решения.
\end{enumerate}

Следовательно, по результатам анализа примеров, разработанная фитнес-функция эффективно оценивает \texttt{SQL}-запросы, направляя генетический алгоритм к решениям, восстанавливающим целостность данных. 

Важными факторами являются выбор релевантных критериев, учет особенностей задачи при их интерпретации, а также правильная настройка параметров генетического алгоритма для достижения оптимального баланса между точностью и эффективностью поиска. Сочетание этих элементов обеспечивает перспективность подхода для автоматической коррекции ошибок в базах данных.

\conclusion
В данной работе была исследована возможность применения генетичес\-кого алгоритма для решения задачи коррекции ошибок в базах данных, обус\-ловленных нарушением ограничений целостности. В результате проведенных исследований был разработана система, включающая в себя инструменты для анализа структуры базы данных, генерации \texttt{SQL}-запросов, оценки их соответст\-вия требованиям целостности, эффективности и сохранности данных. 

Ключевым элементом разработанной системы является специализированная фитнес-функция, учитывающая баланс между соблюдением ограничений, соответствием схеме и минимизацией изменений данных.

Проведенные эксперименты, реализованные на языке \texttt{Java} с использованием библиотеки \texttt{JDBC} для взаимодействия с СУБД и \texttt{Swing} для создания графического интерфейса, позволили оценить эффективность предложенного подхода на тестовой базе данных. Полученные результаты подтвердили перспективность использования генетического алгоритма для решения задачи коррекции ошибок. 


Было показано, что разработанный алгоритм способен находить решения для исправления различных типов ошибок, возникающих вследствие человеческого фактора, программных сбоев или иных причин. Достигнутые результаты свидетельствуют о возможности автоматизации процесса коррекции данных, снижении трудозатрат и повышении качества информационных систем.

Несмотря на полученные положительные результаты, следует отметить некоторые ограничения, связанные с вычислительной сложностью генетичес\-кого алгоритма и его чувствительностью к выбору параметров. Дальнейшие исследования могут быть направлены на решение следующих задач: оптимизация производительности алгоритма за счет применения параллельных вычислений и разработки более эффективных генетических операторов; адаптация фитнес-функции для работы со сложными и специфическими типами ошибок, характерными для различных предметных областей; расширение функциональности программного комплекса за счет интеграции с популярными системами управления базами данных и улучшения пользовательского интерфейса.

Результаты данной работы демонстрируют потенциал генетических алгоритмов для решения задачи коррекции ошибок в базах данных, что может способствовать снижению затрат на поддержку данных, повышению их качества и обеспечению надежности информационных систем. 

Полученные знания и разработанные решения представляют интерес для дальнейших исследований в области интеллектуальной обработки данных и могут быть использованы в практических разработках для автоматизации процессов управления данными.

\printthebibliography
\end{document}

\begin{pmk-biblio}
\bibitem{bib1}
Генетические алгоритмы : учебно-методическое пособие / сост. А. А. Короткин, М. В. Краснов ; Яросл. гос. ун-т им. П. Г. Демидова. –– Ярос\-лавль : ЯрГУ, 2020. ––  Ярославль, 2020. –– 40 с.

\bibitem{bib2}
Генетические алгоритмы — математический аппарат.~ Loginom~ : [сайт]. ~--- URL: \url{https://loginom.ru/blog/ga-math} (дата обращения: 22.05.2024). ~--- Загл. с титул. экрана.

\bibitem{bib3}
Генетическое программирование.~ Интуит~ : [сайт]. ~--- URL: \url{https://intuit.ru/studies/curriculums/19209/courses/1284/lecture/24178} (дата обращения: 22.05.2024). ~--- Загл. с титул. экрана.

\bibitem{bib4}
Дудаков, С. М. Сложность обновлений дедуктивных баз данных с фиксированным ограничением целостности // Вестник ТвГУ. Серия: Прик\-ладная математика.~--- 2003.~--- № 1.~--- С. 16–27. ISSN 1995-0136.

\bibitem{bib5}
Руководство по языку программирования Java.~ Сайт о программировании~ : [сайт]. ~--- URL: \url{https://metanit.com/java/tutorial} (дата обращения: 22.05.2024). ~--- Загл. с титул. экрана.

\bibitem{bib6}
Creating a genetic algorithm for beginners.~ The Project Spot~ : [сайт]. ~--- URL: \url{https://www.theprojectspot.com/tutorial-post/creating-a-genetic-algorithm-for-beginners/3} (дата обращения: 22.05.2024). ~--- Загл. с титул. экрана.

\bibitem{bib7}
Java Swing: знакомство.~ Progoschool~ : [сайт]. ~--- URL: \url{https://progoschool.ru/java/java-swing} (дата обращения: 22.05.2024). ~--- \mbox{Загл.} с титул. экрана.


\bibitem{bib8}
PostgreSQL 17.5 Documentation.~ The PostgreSQL Global Development~ : [сайт]. ~--- URL: \url{https://www.postgresql.org/docs/17/index.html} (дата обращения: 22.05.2024). ~--- Загл. с титул. экрана.
\end{pmk-biblio}


\end{document}
